% Noether P08 Korean producer draft — UNCHECKED
% T06-U22; ED0002 lines6183--6187; source 672 LF B / 80F180E0ADABC0A9972DC61BE5CFE94675B91579FBB261F1447CC7F64A1250A5
% Pointer v009 / ED0002: B06BE3530D9CF2E82B56FDBA7FE41D5D044DF2425DFA2A059D4939EAA2F7A6C2 / C9A125167ACB33D914EE4374B65AE7CDF0052F568371B8B77B720EA178ABF0E3
% Translation only; no review, build, render, or approval.

그러므로 a)와 b)에 의해 $\rho=\tau$이다. 또한 $\mathcal T$는 $\mathcal L$의 부분족이므로, 이로써 두 선형 형식족의 동치가 증명되었다. 따라서 $\mathcal L$의 모든 형식, 특히 $\theta$도 선형독립인 $\rho$개의 형식 $h(x,y,z)$의 유리수 계수 선형결합으로 나타낼 수 있다.
\[
  \theta=\sum c_i h^{(i)}(x,y,z)=\sum PZ.
\]
$\theta$의 부정계수 $a$에 대해 유도한 이 항등식은 부정원 $a$를 임의의 유리성 영역의 원소들로 바꾸어도 그대로 성립한다. 따라서 \emph{세 개의 이진 변수열인 경우에 환원정리가 일반적으로 증명되었다.}
